% Auto-generated by 03_analysis/83_preclosure_timing_tests.R -- do not edit by hand.
\begin{table}[!htbp]\centering
\caption{Pre-closure timing tests: is closure timing predicted by pre-existing local conditions or trends?}
\label{tab:preclosure-timing-tests}
\begin{threeparttable}
\small
\begin{tabular}{lccc}
\toprule
Standardized predictor & Coef. & Std.\ err. & $p$-value \\
\midrule
\multicolumn{4}{l}{\emph{Panel A. One-predictor-at-a-time OLS, DV $=$ closure year}}\\
\addlinespace
Pre-closure suicide trend (catchment) & -0.145 & 0.195 & 0.464 \\
Pre-closure self-harm trend (catchment) & -0.387** & 0.185 & 0.044 \\
Pre-closure psych-admissions trend & -0.386** & 0.185 & 0.045 \\
Pre-closure total-admissions trend & -- & -- & -- \\
Pre-closure travel-burden trend & 0.073 & 0.196 & 0.710 \\
Pre-closure admission decline (\%) & 0.280 & 0.239 & 0.247 \\
Baseline suicide rate (catchment) & -0.034 & 0.168 & 0.847 \\
Baseline self-harm rate (catchment) & -0.030 & 0.169 & 0.863 \\
Baseline psychiatric admissions & 0.163 & 0.241 & 0.504 \\
Baseline CAPS facilities & 0.125 & 0.242 & 0.608 \\
Baseline log population & 0.153 & 0.115 & 0.208 \\
Baseline log GDP & 0.149 & 0.115 & 0.221 \\
\addlinespace
\midrule
\multicolumn{4}{l}{\emph{Panel B. Parsimonious joint OLS, DV $=$ closure year}}\\
\addlinespace
Pre-closure admission decline (\%) & 0.507* & 0.282 & 0.078 \\
Baseline psychiatric admissions & 0.479 & 0.287 & 0.102 \\
Baseline CAPS facilities & 0.155 & 0.247 & 0.533 \\
\addlinespace
\midrule
\multicolumn{4}{l}{\emph{Panel C. Pre-trends-only joint OLS, DV $=$ closure year}}\\
\addlinespace
Pre-closure suicide trend (catchment) & -0.015 & 0.198 & 0.940 \\
Pre-closure self-harm trend (catchment) & -0.392* & 0.201 & 0.060 \\
Pre-closure psych-admissions trend & -0.419** & 0.177 & 0.025 \\
Pre-closure travel-burden trend & 0.130 & 0.180 & 0.476 \\
\addlinespace
\midrule
\multicolumn{4}{l}{\textbf{Joint $F$-test, Panel~B (closure year): $F(3,44)=1.43$, $p=0.248$; adj.\ $R^2=0.0264$; $N=48$.}}\\
\multicolumn{4}{l}{\textbf{Joint $F$-test, Panel~C pre-trends (closure year): $F(4,29)=2.39$, $p=0.074$; adj.\ $R^2=0.1438$; $N=34$.}}\\
\multicolumn{4}{l}{\textbf{Joint test (early vs.\ late closure): LR $\chi^2=4.91$, $p=0.178$; $N=48$.}}\\
\multicolumn{4}{l}{Joint test (PNASH exact $\pm1$ window, LPM): $F=1.58$, $p=0.207$; $N=48$.}\\
\bottomrule
\end{tabular}
\begin{tablenotes}[flushleft]\footnotesize
\item \textit{Notes.} Unit of observation is the closure event ($N=48$). All predictors are standardized to mean~0 and standard deviation~1, so coefficients are directly comparable in standard-deviation units; a positive coefficient on the closure-year DV means the predictor is associated with a \emph{later} closure. Catchment pre-trends are per-closure linear slopes (population-weighted catchment mean regressed on year) over the pre-period years $t<g$; closures in the earliest cohort have only two pre-years and so receive a missing slope (reported as `--' / dropped from the joint model). The total-admissions trend is not reported because the municipality--year panel carries no total-admissions series. Baseline catchment suicide/self-harm rates, population, and GDP are available for a subset of events; Panel~B is therefore restricted to the fully covered standardized scalars (top predictors by univariate $|t|$, capped to avoid overfitting at $N=48$) and Panel~C reports a complementary joint specification using only the catchment pre-trends, which directly tests whether timing is driven by pre-existing local trends. The early-vs-late dummy equals one if the closure year is below the sample median (2014); the PNASH exact $\pm1$ indicator (37 within / 11 outside) is modeled by a linear probability model because the logit suffers perfect separation. ${}^{*}p<0.10$, ${}^{**}p<0.05$, ${}^{***}p<0.01$.
\item \textit{Power caveat.} With $N=48$ these tests are low-powered. A failure to reject predictability of timing is \emph{consistent with} as-good-as-random closure timing given the PNASH inspection-cycle institutional setting, but it is not proof of unconfoundedness.
\end{tablenotes}
\end{threeparttable}
\end{table}
